\clearpage
\section{Recurrent language models}\label{sec:rnn}
\subsection*{Carry a state instead of fixing a window}
A fixed-window neural LM cannot distinguish prefixes whose final $m$
tokens agree. To let earlier tokens influence a prediction, a recurrent
network carries a state forward. At each step it combines the current
embedding with the previous state:
\[
 e_t=E[x_t],\qquad
 h_t=\tanh(W_xe_t+W_hh_{t-1}+b_h),\qquad
 z_t=W_oh_t+b_o.
\]
The state has width $H$, so $W_x\in\R^{H\times d}$,
$W_h\in\R^{H\times H}$, and $W_o\in\R^{K\times H}$.
The output $\softmax(z_t)$ predicts $x_{t+1}$. The same parameters are
used at every position, while the state changes with the sequence.

\fig{recurrent-state}{Unrolling a recurrent network exposes the sequence of
computations. Repeated boxes share parameters; the hidden states are different tensors.}

There is no fixed architectural cutoff at $m$ previous tokens in this
recurrence. An earlier input can affect later states through repeated
updates. This describes a possible dependency, not a guarantee that useful
information will be retained or learned. A finite-width, finite-precision
state must preserve relevant distinctions through those updates.

\subsection*{Worked scalar recurrence}
To isolate the state update, choose $h_0=0$, $W_x=1$, $W_h=0.5$,
$b_h=0$, and scalar inputs $(e_1,e_2,e_3)=(1,0,0)$. Then
\[
 h_1=\tanh(1)\approx0.7616,\quad
 h_2=\tanh(0.5h_1)\approx0.3634,\quad
 h_3=\tanh(0.5h_2)\approx0.1797.
\]
The first input still affects the third state, although its effect has
shrunk. These are hand-chosen parameters illustrating a mechanism, not
the output of a trained language model.

At a new independent document, reset or explicitly initialize the state.
Carrying a previous document's state would change the conditioning context.
For padded batches, exclude padding positions from the loss unless they
are intentionally part of the modeled event space. Reset policy and loss
masking are data conventions, not details to leave implicit.

\takeaway{Recurrence shares a transition rule across time. It trades a
fixed input window for a repeatedly updated summary of the prefix.}
\selfcheck{Does making a batch larger remove the dependence of $h_t$ on
$h_{t-1}$ within one sequence? No; batch parallelism and time dependence are separate.}
\textbf{Reading connection.} \readingref{R08} develops learned temporal
representations; \readingref{R11} applies recurrent states to language modeling.

\nextpage{Backpropagation through time}
Training an RNN differentiates the unrolled computation graph. A parameter
can affect the loss directly at the current step and indirectly through
many later states. Backpropagation through time (BPTT) applies the ordinary
chain rule to all these paths, then sums contributions to each shared parameter.

\fig{bptt-paths}{Forward states connect adjacent time steps. Backward signals
from a later loss pass through earlier states. A truncation boundary blocks
the backward connection while a numerical state can still be carried forward.}

For the vanilla recurrence, hold input embeddings fixed and define the
local recurrent Jacobian
\[
 J_t=\frac{\partial h_t}{\partial h_{t-1}}
 =\operatorname{diag}(1-h_t\odot h_t)W_h.
\]
The influence of an earlier state on a later one is a product:
\[
 \frac{\partial h_t}{\partial h_s}=J_tJ_{t-1}\cdots J_{s+1},\qquad s<t.
\]
The order matters because matrices do not generally commute. A loss at
step $t$ sends its derivative through this product to states before $t$.
When differentiating a shared weight, sum over the steps where that weight
was used as well as over the losses affected by each use.

\textbf{A small graph to trace.} With three recurrent steps and losses at
all three outputs, the first state receives contributions from all three
losses. The third state receives a contribution only from the third loss
in this graph. The transition weights are shared, so their gradient
accumulates contributions at steps one, two, and three.

\subsection*{Truncation changes the learning path}
For long streams, truncated BPTT divides computation into chunks. The
final state of one chunk may initialize the next, while its graph is
detached at the boundary. The next chunk sees earlier information in the
forward state, but its loss cannot change parameters through the detached
earlier computation. The forward context and backward credit-assignment
horizon are therefore different quantities.

For a first-order RNN, ordinary full BPTT stores information proportional
to sequence length for the backward pass. Checkpointing can trade extra
computation for less storage, but it does not alter the mathematical
gradient when the recomputation is exact. Truncation instead omits paths
and generally changes that gradient.

\selfcheck{If a clue is carried across a detached boundary, can the current
loss teach the earlier transition how to store that clue through that path?
No. It can still update parameters through their uses in the current chunk.}

\nextpage{Why long dependencies are difficult to learn}
Repeated multiplication can make sensitivity very small or very large.
For compatible matrix norms,
\[
 \left\|\frac{\partial h_t}{\partial h_s}\right\|
 \le\prod_{k=s+1}^{t}\|J_k\|.
\]
If every factor is bounded by a constant below one, the bound decreases
exponentially with distance. Factors larger than one permit amplification,
but do not by themselves prove explosion in every direction. Singular
directions, nonlinearities, and the actual state trajectory matter.

\fig{recurrent-gradients}{Exact scalar powers illustrate repeated multiplication.
The curves are not measured RNN gradients: real gradients involve matrix
products and contributions from several paths.}

At 50 steps, $0.9^{50}\approx0.00515$ and $1.1^{50}\approx117.39$.
The scalar example makes the issue visible without assuming that an RNN's
Jacobian is constant. A saturated $\tanh$ also has a small derivative,
so a model can represent a dependency in principle while receiving a weak
learning signal for it. \readingref{R09} studies the difficulty of learning
long dependencies with gradient descent.

\subsection*{What gradient clipping does}
For a gradient vector $g$ and threshold $\tau>0$, global norm clipping uses
\[
 \widetilde g=
 \begin{cases}
 g,&\|g\|_2\le\tau,\\
 \tau g/\|g\|_2,&\|g\|_2>\tau.
 \end{cases}
\]
For $g=(3,4)$ and $\tau=2$, the clipped vector is $(1.2,1.6)$.
It preserves direction and caps the norm of the gradient supplied to the
optimizer. With plain SGD it also bounds that update norm by $\eta\tau$;
adaptive optimizers can rescale coordinates afterward. Clipping does not
restore a learning signal that has already vanished.

\takeaway{Long-range learning depends on both forward memory and backward
sensitivity. Neither a long theoretical context nor gradient clipping
guarantees successful credit assignment.}
\selfcheck{Would setting every recurrent weight to one guarantee stable gradients?
No: the relevant factors are full local Jacobians, including nonlinear derivatives.}
\textbf{Supporting reading.} \href{https://proceedings.mlr.press/v28/pascanu13.html}{Pascanu,
Mikolov, and Bengio (2013)} analyzes these training difficulties and clipping.
